\documentclass[10pt, letterpaper]{article}
\usepackage[utf8]{inputenc}
\usepackage[T1]{fontenc}
\usepackage{geometry}
\usepackage{enumitem}
\usepackage{hyperref}
\usepackage{titlesec}
\usepackage{color}
\usepackage{fancyhdr}

% Margins
\geometry{top=0.75in, bottom=0.75in, left=1in, right=1in}

% Colors
\definecolor{columbiablue_official}{RGB}{0, 33, 71}

% Header
\pagestyle{fancy}
\fancyhf{}
\renewcommand{\headrulewidth}{0.5pt}
\renewcommand{\headrule}{\hbox to\headwidth{\color{columbiablue_official}\leaders\hrule height \headrulewidth\hfill}}
\lhead{\small \textbf{\color{columbiablue_official} COLUMBIA UNIVERSITY} \\ \href{https://www.arch.columbia.edu/}{Graduate School of Architecture, Planning and Preservation}}
\rhead{\small Semi-Structured Stakeholder Interview Guide}

% Section Styling
\titleformat{\section}{\large\bfseries\color{columbiablue_official}}{}{0em}{}[\color{columbiablue_official}\titlerule]
\titlespacing*{\section}{0pt}{12pt}{6pt}

\begin{document}

\begin{center}
    {\Large \textbf{\color{columbiablue_official} Predicting NIMBYism in Austin’s Housing Reform Era}}
\end{center}

\section*{Purpose}
To explore how key stakeholders understand neighborhood opposition to housing development in Austin and how they view the potential use of predictive or algorithmic tools in planning and zoning.

\section*{Opening}
\begin{enumerate}[leftmargin=*, label=\arabic*.]
    \item Thank the participant for their time; confirm they have read and signed the consent form.
    \item Confirm whether they consent to audio recording.
    \item Remind them they may skip any question or stop at any time.
\end{enumerate}

\section*{Part 1: Background and Role}
\begin{enumerate}[leftmargin=*, label=\arabic*.]
    \item Can you briefly describe your current role and how it relates to housing, planning, or land use in Austin?
    \item How long have you been involved with housing or land use issues in Austin?
    \item In your role, what kinds of development or zoning decisions do you most often encounter?
\end{enumerate}

\section*{Part 2: Experiences with Neighborhood Opposition}
\begin{enumerate}[leftmargin=*, label=\arabic*.]
    \item When you think of “neighborhood opposition” to housing or development projects in Austin, what comes to mind?
    \item Can you describe a recent case or example where neighborhood opposition played a major role? What happened?
    \item In your experience, what types of projects are most likely to generate opposition?
    \item Which kinds of people or groups tend to be most active in opposing projects?
\end{enumerate}

\section*{Part 3: Perceptions of Fairness and Participation}
\begin{enumerate}[leftmargin=*, label=\arabic*.]
    \item Do you think current processes for voicing opposition (such as public hearings, petitions, and appeals) are fair? Why or why not?
    \item Are there groups you feel are over-represented or under-represented in these processes?
    \item How do you think opposition affects whose interests are ultimately reflected in housing and zoning decisions?
\end{enumerate}

\section*{Part 4: Views on Predictive and Algorithmic Tools}
\begin{enumerate}[leftmargin=*, label=\arabic*.]
    \item Some cities are considering or experimenting with tools that use data to predict where opposition to development is likely to occur. Have you encountered ideas like this before?
    \item How do you feel about the idea of a model that forecasts which areas or projects will face higher levels of opposition?
    \item What potential benefits do you see, if any, in using such tools for planning or outreach?
    \item What potential risks or harms concern you most? For example, concerns about profiling, surveillance, or chilling participation.
    \item Under what conditions, if any, would you consider the use of such tools acceptable or legitimate? (e.g., transparency, public input, limits on use).
\end{enumerate}

\section*{Part 5: Alternative Technological Approaches}
\begin{enumerate}[leftmargin=*, label=\arabic*.]
    \item Instead of predicting opposition, are there other ways you think data or technology could improve planning and housing processes in Austin?
    \item Are there tools or approaches used in other cities that you find promising?
    \item If you could design one tool or system to improve public engagement or make development decisions fairer, what would it look like?
\end{enumerate}

\section*{Closing}
\begin{enumerate}[leftmargin=*, label=\arabic*.]
    \item Is there anything else you would like to add about neighborhood opposition, planning, or the use of data and technology in housing decisions?
    \item Thank the participant again and remind them whom to contact if they have further questions.
\end{enumerate}

\vfill
\noindent \textbf{IRB Contact at Columbia (Morningside)} \\
If you have questions about your rights or welfare as a research participant, you may contact: \\
Columbia University Morningside Institutional Review Board (IRB) \\
Phone: (212) 851-7040 \\
Email: \href{mailto:AskIRB@columbia.edu}{AskIRB@columbia.edu}

\end{document}
